\documentclass[11pt]{article}

\usepackage[T1]{fontenc}
\usepackage{lmodern}
\usepackage{microtype}
\usepackage[margin=0.82in,headheight=15pt]{geometry}
\usepackage{mathtools,amssymb,amsthm}
\usepackage{xcolor}
\usepackage[most]{tcolorbox}
\usepackage{enumitem}
\usepackage{fancyhdr}
\usepackage{hyperref}

\definecolor{ink}{HTML}{162332}
\definecolor{navy}{HTML}{17324D}
\definecolor{teal}{HTML}{167D82}
\definecolor{gold}{HTML}{D69E2E}
\definecolor{coral}{HTML}{D8664A}
\definecolor{mist}{HTML}{EDF5F4}
\definecolor{paper}{HTML}{FBFAF7}
\definecolor{muted}{HTML}{5D6B78}

\pagecolor{paper}
\color{ink}
\hypersetup{
  colorlinks=true,
  linkcolor=teal,
  urlcolor=coral,
  pdfauthor={Igor Rivin},
  pdftitle={Number Theory Problem Set 1: Euclid and the Architecture of the Integers}
}

\setlength{\parindent}{0pt}
\setlength{\parskip}{4pt}
\setlist[enumerate,1]{label=\textbf{\arabic*.},leftmargin=*,itemsep=10pt,topsep=6pt}
\setlist[enumerate,2]{label=(\alph*),leftmargin=1.75em,itemsep=4pt,topsep=4pt}

\newcommand{\Z}{\mathbb Z}
\newcommand{\N}{\mathbb N}
\newcommand{\divides}{\mid}
\newcommand{\ndivides}{\nmid}
\newcommand{\points}[1]{\hfill\textcolor{coral}{\small\bfseries[#1 points]}}
\newcommand{\edition}{\ifsolutions Worked Solutions\else Student Handout\fi}

\pagestyle{fancy}
\fancyhf{}
\fancyhead[L]{\small\color{muted}Problem Set 1}
\fancyhead[R]{\small\color{muted}Patterns, Proofs, and Primes}
\fancyfoot[C]{\color{teal}\thepage}
\renewcommand{\headrulewidth}{0.3pt}

\newtcolorbox{overviewbox}{
  enhanced,breakable,colback=mist,colframe=teal,boxrule=.8pt,arc=2mm,
  left=4mm,right=4mm,top=2.5mm,bottom=2.5mm,
  borderline west={3pt}{0pt}{teal}
}

\newtcolorbox{instructornote}{
  enhanced,breakable,colback=gold!7,colframe=gold,boxrule=.7pt,arc=2mm,
  title={Instructor note},fonttitle=\bfseries\color{navy}
}

\NewDocumentEnvironment{solution}{+b}{%
  \ifsolutions
    \begin{tcolorbox}[
      enhanced,breakable,colback=white,colframe=teal!55,boxrule=.55pt,arc=1mm,
      left=3mm,right=3mm,top=2mm,bottom=2mm,
      title={\small\bfseries Solution},title after break={\small\bfseries Solution (continued)},
      fonttitle=\small\bfseries,coltitle=navy]
      \small #1
    \end{tcolorbox}
  \fi
}{}

\begin{document}

\begin{tcolorbox}[
  enhanced,colback=paper,colframe=navy,boxrule=1.1pt,arc=2mm,
  left=5mm,right=5mm,top=3.5mm,bottom=3.5mm,
  borderline west={5pt}{0pt}{teal}]
  {\Large\bfseries\color{navy}Problem Set 1: Euclid and the Architecture of the Integers}
  \hfill{\bfseries\color{teal}\edition}\\[3pt]
  {\large Divisibility, gcds, primes, and unique factorization}\\[4pt]
  {\small\color{muted}80 points + 10 optional challenge points}
  \ifsolutions\else
    \\[7pt]\textbf{Name:}\ \rule{0.46\linewidth}{0.4pt}\hfill
    \textbf{Due:}\ \rule{0.21\linewidth}{0.4pt}
  \fi
\end{tcolorbox}

\begin{overviewbox}
\textbf{Purpose.} This assignment develops the two complementary habits of
elementary number theory: compute examples efficiently, then explain why the
method works. Show the divisions in every use of Euclid's algorithm and name
the theorem used in every proof.

\textbf{You may use:} the division algorithm, the Euclidean algorithm,
B\'ezout's identity, Euclid's lemma, and the fundamental theorem of arithmetic.
Unless a problem explicitly asks for a decimal answer, exact arithmetic is expected.
\end{overviewbox}

\ifsolutions
\begin{instructornote}
Problem 3 is deliberately removable. If B\'ezout's identity has not yet been
covered, assign 3(a) as a guided extension and postpone 3(b)--(c). The remaining
problems do not depend on Problem 3.
\end{instructornote}
\fi

\section*{I. Divisibility and Euclid's algorithm}

\begin{enumerate}
\item \textbf{The invariant behind the algorithm.} \points{6}
Let $a,b,d,m,n\in\Z$.
\begin{enumerate}
  \item Prove that if $d\divides a$ and $d\divides b$, then
  $d\divides(ma+nb)$.
  \item Suppose $a=qb+r$. Prove that $(a,b)$ and $(b,r)$ have exactly the
  same common divisors. Conclude that $\gcd(a,b)=\gcd(b,r)$.
\end{enumerate}
\begin{solution}
\begin{enumerate}
  \item Since $d\divides a$ and $d\divides b$, write $a=du$ and $b=dv$ for
  some $u,v\in\Z$. Then
  \[
    ma+nb=mdu+ndv=d(mu+nv),
  \]
  so $d\divides(ma+nb)$.
  \item If $d$ divides both $a$ and $b$, part (a), with coefficients $1$ and
  $-q$, shows that $d\divides a-qb=r$. Thus $d$ divides both $b$ and $r$.
  Conversely, if $d\divides b$ and $d\divides r$, then
  $d\divides qb+r=a$. The pairs therefore have the same common divisors, so
  their greatest positive common divisors are equal.
\end{enumerate}
\end{solution}

\item \textbf{Three runs of Euclid's algorithm.} \points{9}
Compute each gcd, displaying every division. For part (c), describe the pattern
in the remainders.
\begin{enumerate}
  \item $\gcd(2024,748)$
  \item $\gcd(12345,6789)$
  \item $\gcd(4181,2584)$
\end{enumerate}
\begin{solution}
\begin{enumerate}
  \item
  \begin{align*}
  2024&=2\cdot748+528,&748&=1\cdot528+220,\\
  528&=2\cdot220+88,&220&=2\cdot88+44,\\
  88&=2\cdot44.
  \end{align*}
  Hence $\gcd(2024,748)=44$.

  \item
  \begin{align*}
  12345&=1\cdot6789+5556,&6789&=1\cdot5556+1233,\\
  5556&=4\cdot1233+624,&1233&=1\cdot624+609,\\
  624&=1\cdot609+15,&609&=40\cdot15+9,\\
  15&=1\cdot9+6,&9&=1\cdot6+3,\\
  6&=2\cdot3.
  \end{align*}
  Thus $\gcd(12345,6789)=3$.

  \item
  \begin{align*}
  4181&=1\cdot2584+1597,&2584&=1\cdot1597+987,\\
  1597&=1\cdot987+610,&987&=1\cdot610+377,\\
  610&=1\cdot377+233,&377&=1\cdot233+144,\\
  233&=1\cdot144+89,&144&=1\cdot89+55,\\
  89&=1\cdot55+34,&55&=1\cdot34+21,\\
  34&=1\cdot21+13,&21&=1\cdot13+8,\\
  13&=1\cdot8+5,&8&=1\cdot5+3,\\
  5&=1\cdot3+2,&3&=1\cdot2+1,\\
  2&=2\cdot1.
  \end{align*}
  Therefore $\gcd(4181,2584)=1$. The remainders run backward through the
  Fibonacci numbers; each quotient before the last is $1$.
\end{enumerate}
\end{solution}

\item \textbf{Running Euclid backward.} \points{12}
\begin{enumerate}
  \item Use the divisions from Problem 2(a) to express $44$ as
  $2024x+748y$ with $x,y\in\Z$.
  \item Find one integer solution of $2024x+748y=132$, and then find all
  integer solutions.
  \item Explain why $2024x+748y=100$ has no integer solution.
\end{enumerate}
\begin{solution}
\begin{enumerate}
  \item Reverse the divisions:
  \begin{align*}
  44&=220-2\cdot88\\
    &=220-2(528-2\cdot220)=5\cdot220-2\cdot528\\
    &=5(748-528)-2\cdot528=5\cdot748-7\cdot528\\
    &=5\cdot748-7(2024-2\cdot748)
      =-7\cdot2024+19\cdot748.
  \end{align*}
  Thus $(x,y)=(-7,19)$.
  \item Multiplying the identity in (a) by $3$ gives one solution
  $(x_0,y_0)=(-21,57)$. Since $2024/44=46$ and $748/44=17$, all solutions are
  \[
    x=-21+17t,\qquad y=57-46t,\qquad t\in\Z.
  \]
  Indeed, the changes contribute $2024(17)-748(46)=0$, and every solution
  differs from the displayed one by such a change.
  \item Every integer linear combination of $2024$ and $748$ is divisible by
  their gcd, $44$. Since $44\ndivides100$, no solution exists.
\end{enumerate}
\end{solution}

\item \textbf{Useful gcd identities.} \points{8}
\begin{enumerate}
  \item Prove that $\gcd(a+kb,b)=\gcd(a,b)$ for every $k\in\Z$.
  \item Use part (a), without factoring, to prove that
  $\gcd(2n+1,n)=1$ for every positive integer $n$.
  \item Prove that two consecutive positive integers are relatively prime.
\end{enumerate}
\begin{solution}
\begin{enumerate}
  \item A common divisor of $a$ and $b$ divides $(a+kb)$ and $b$. Conversely,
  a common divisor of $a+kb$ and $b$ divides $(a+kb)-kb=a$. The pairs have
  exactly the same common divisors.
  \item Apply (a) with $a=1$, $b=n$, and $k=2$:
  $\gcd(2n+1,n)=\gcd(1,n)=1$.
  \item A common divisor of $n$ and $n+1$ divides their difference $1$.
  Hence their gcd is $1$.
\end{enumerate}
\end{solution}

\item \textbf{Why the algorithm stops and succeeds.} \points{7}
In a paragraph, prove the Euclidean algorithm: explain both why repeated
division must terminate and why the last nonzero remainder is the gcd of the
original two positive integers.
\begin{solution}
At each step the new remainder is a nonnegative integer strictly smaller than
the preceding positive remainder. There cannot be an infinite strictly
decreasing sequence of nonnegative integers, so eventually a remainder is
zero. By Problem 1(b), replacing a pair by the divisor and remainder does not
change its common divisors. Thus every pair in the algorithm has the same gcd
as the original pair. At the final nonzero step we have $r_{k-1}=q r_k$, so
the common divisors of $r_{k-1}$ and $r_k$ are precisely the divisors of $r_k$.
Their greatest positive common divisor is therefore $r_k$, the last nonzero
remainder.
\end{solution}
\end{enumerate}

\section*{II. Primes and unique factorization}

\begin{enumerate}[resume]
\item \textbf{Prime-exponent coordinates.} \points{12}
\begin{enumerate}
  \item Find the prime factorizations of $7560$ and $8820$.
  \item Use the exponent vectors, without applying Euclid's algorithm, to find
  their gcd and lcm.
  \item How many positive divisors does $7560$ have? Justify the counting rule
  you use.
  \item Verify numerically that
  $\gcd(a,b)\operatorname{lcm}(a,b)=ab$ for this pair.
\end{enumerate}
\begin{solution}
\begin{enumerate}
  \item
  \[
    7560=2^3\cdot3^3\cdot5\cdot7,
    \qquad
    8820=2^2\cdot3^2\cdot5\cdot7^2.
  \]
  \item Coordinatewise minima and maxima give
  \[
  \gcd(7560,8820)=2^2\cdot3^2\cdot5\cdot7=1260
  \]
  and
  \[
  \operatorname{lcm}(7560,8820)=2^3\cdot3^3\cdot5\cdot7^2=52920.
  \]
  \item A positive divisor has the form $2^a3^b5^c7^d$ with
  $0\le a\le3$, $0\le b\le3$, and $0\le c,d\le1$. The choices are
  independent, so there are
  $(3+1)(3+1)(1+1)(1+1)=64$ positive divisors.
  \item
  \[
    1260\cdot52920=66{,}679{,}200
    =7560\cdot8820.
  \]
\end{enumerate}
\end{solution}

\item \textbf{Euclid's lemma, the hinge of uniqueness.} \points{8}
\begin{enumerate}
  \item Let $p$ be prime. Prove that if $p\divides ab$, then
  $p\divides a$ or $p\divides b$. You may use B\'ezout's identity.
  \item Deduce that if $p\divides a^m$ for an integer $m\ge1$, then
  $p\divides a$.
  \item In two or three sentences, explain where part (a) enters the uniqueness
  proof in the fundamental theorem of arithmetic.
\end{enumerate}
\begin{solution}
\begin{enumerate}
  \item If $p\divides a$, there is nothing to prove. Otherwise
  $\gcd(p,a)=1$, since the only positive divisors of $p$ are $1$ and $p$.
  B\'ezout gives $px+ay=1$. Multiplying by $b$ yields
  $pbx+aby=b$. Both terms on the left are divisible by $p$ (the second because
  $p\divides ab$), so $p\divides b$.
  \item Apply part (a) to $a^m=a\cdot a^{m-1}$. Either $p\divides a$, or
  $p\divides a^{m-1}$. Repeating this argument eventually gives
  $p\divides a$.
  \item If $p_1\cdots p_r=q_1\cdots q_s$, then $p_1$ divides the product on
  the right. Euclid's lemma implies that $p_1$ equals one of the prime factors
  $q_j$. Cancel that common factor and repeat; this forces the two lists of
  primes to agree up to order.
\end{enumerate}
\end{solution}

\item \textbf{Two consequences of prime factorization.} \points{8}
\begin{enumerate}
  \item Suppose $\gcd(a,b)=1$ and $ab$ is a perfect square. Prove that both
  $a$ and $b$ are perfect squares.
  \item Prove that $\sqrt p$ is irrational for every prime $p$.
\end{enumerate}
\begin{solution}
\begin{enumerate}
  \item In the prime factorizations of $a$ and $b$, no prime occurs in both,
  because $\gcd(a,b)=1$. Since $ab$ is a square, every prime has even exponent
  in the factorization of the product. Each exponent comes entirely from
  either $a$ or $b$, so every exponent in each factor is even. Hence $a$ and
  $b$ are squares.
  \item Suppose $\sqrt p=m/n$ with positive integers $m,n$ chosen relatively
  prime. Then $m^2=pn^2$, so $p\divides m^2$. Problem 7(b) gives
  $p\divides m$; write $m=pk$. Substitution gives $p^2k^2=pn^2$, hence
  $n^2=pk^2$ and $p\divides n^2$. Again $p\divides n$, contradicting
  $\gcd(m,n)=1$. Therefore $\sqrt p$ is irrational.
\end{enumerate}
\end{solution}

\item \textbf{There is no last prime.} \points{10}
\begin{enumerate}
  \item Give Euclid's proof that there are infinitely many primes. Your proof
  must explicitly explain why the constructed integer has a prime divisor and
  why that divisor is new.
  \item The first six primes are $2,3,5,7,11,13$. Compute their product plus
  one and verify
  \[
    2\cdot3\cdot5\cdot7\cdot11\cdot13+1=30031=59\cdot509.
  \]
  What common misunderstanding of Euclid's proof does this example correct?
\end{enumerate}
\begin{solution}
\begin{enumerate}
  \item Suppose there were only finitely many primes, say
  $p_1,\ldots,p_k$. Form $N=p_1p_2\cdots p_k+1$. Since $N>1$, the existence
  part of the fundamental theorem of arithmetic says that $N$ has a prime
  divisor $q$. But division of $N$ by any $p_i$ leaves remainder $1$, so none
  of the listed primes divides $N$. Thus $q$ is a prime missing from the
  supposedly complete list, a contradiction.
  \item The product is $30030$, so adding $1$ gives $30031$. Direct
  multiplication verifies
  $59\cdot509=59(500+9)=29500+531=30031$.
  This shows that the number constructed in Euclid's proof need not itself be
  prime. The proof needs only a prime divisor of that number which is absent
  from the proposed finite list.
\end{enumerate}
\end{solution}
\end{enumerate}

\section*{Optional challenge}

\begin{enumerate}[resume]
\item \textbf{Infinitely many primes leaving remainder $3$ upon division by $4$.}
\points{10 extra}
Suppose, for contradiction, that the complete list of such primes is
$p_1,\ldots,p_k$, and consider
\[
  N=4p_1p_2\cdots p_k-1.
\]
Prove that some prime divisor of $N$ leaves remainder $3$ upon division by $4$,
and that this divisor is not one of the $p_i$. Conclude that there are
infinitely many primes of the form $4n+3$.
\begin{solution}
The integer $N$ is odd and leaves remainder $3$ upon division by $4$. Every odd
prime leaves remainder $1$ or $3$ upon division by $4$. If every prime factor
of $N$ left remainder $1$, then their product (with multiplicity) would also
leave remainder $1$, contradicting the fact that $N$ leaves remainder $3$.
Thus at least one prime divisor $q$ of $N$ leaves remainder $3$. For each
$p_i$, division of $N$ by $p_i$ leaves remainder $-1$, so $p_i\ndivides N$.
Therefore $q$ is a new prime of the form $4n+3$, contradicting completeness of
the list.
\end{solution}
\end{enumerate}

\vfill
\begin{center}
\color{muted}\small
The arithmetic is the evidence; the proof is the explanation.
\end{center}

\end{document}
